%% Delete everything above this line and pdflatex the resulting file twice. %%
%% LaTeX 2e

\pdfoutput=1 % tell arXiv to use PDFLaTeX

\RequirePackage[l2tabu, orthodox]{nag}

\documentclass[12pt,reqno]{amsart}
\usepackage[margin=1in]{geometry} 
\usepackage{url,amssymb,stmaryrd,enumerate,colonequals,graphicx}
\usepackage[all]{xy} % for commutative diagrams
\usepackage{mathrsfs} % for \mathscr (script letters)
\usepackage{comment} % for commenting out large chunks of LaTeX

\usepackage[T1]{fontenc}
\usepackage[utf8]{inputenc}

% Color comments!
\usepackage[dvipsnames,xcdraw]{xcolor}
\newcommand{\bjorn}[1]{{\color{blue} \textsf{$\clubsuit\clubsuit\clubsuit$ Bjorn: [#1]}}}
\newcommand{\blake}[1]{{\color{Mahogany} \textsf{$\triangle$ Blake: [#1]}}}
\newcommand{\rachel}[1]{{\color{red} \textsf{$\diamondsuit$ Rachel: [#1]}}}
\newcommand{\khl}[1]{{\color{magenta} \textsf{KH: [#1]}}}
\newcommand{\coco}[1]{{\color{teal} \textsf{Coco: [#1]}}}
\newcommand{\Shaowu}[1]{{\color{orange} \textsf{$\bigstar$ Shaowu: [#1]}}}

\newcommand{\deficolor}{ForestGreen}
\newcommand{\defi}[1]{\textbf{\color{\deficolor} #1}} % for defined terms
\newcommand{\defimath}[1]{{\color{\deficolor} \text{\boldmath{$#1$}} }} % for defined terms in math mode

\newcommand{\id}{\operatorname{id}}

% Characters
\newcommand{\Aff}{\mathbb{A}}
\newcommand{\C}{\mathbb{C}}
\newcommand{\F}{\mathbb{F}}
\newcommand{\G}{\mathbb{G}}
\newcommand{\bbH}{\mathbb{H}}
\newcommand{\PP}{\mathbb{P}}
\newcommand{\Q}{\mathbb{Q}}
\newcommand{\R}{\mathbb{R}}
\newcommand{\Z}{\mathbb{Z}}
\newcommand{\Qbar}{{\overline{\Q}}}
\newcommand{\Zhat}{{\hat{\Z}}}
\newcommand{\Zbar}{{\overline{\Z}}}
\newcommand{\kbar}{{\overline{k}}}
\newcommand{\Kbar}{{\overline{K}}}
\newcommand{\ksep}{{k^{\operatorname{sep}}}}
\newcommand{\Fbar}{{\overline{\F}}}

% mathcal characters
\newcommand{\calA}{\mathcal{A}}
\newcommand{\calB}{\mathcal{B}}
\newcommand{\calC}{\mathcal{C}}
\newcommand{\calD}{\mathcal{D}}
\newcommand{\calE}{\mathcal{E}}
\newcommand{\calF}{\mathcal{F}}
\newcommand{\calG}{\mathcal{G}}
\newcommand{\calH}{\mathcal{H}}
\newcommand{\calI}{\mathcal{I}}
\newcommand{\calJ}{\mathcal{J}}
\newcommand{\calK}{\mathcal{K}}
\newcommand{\calL}{\mathcal{L}}
\newcommand{\calM}{\mathcal{M}}
\newcommand{\calN}{\mathcal{N}}
\newcommand{\calO}{\mathcal{O}}
\newcommand{\calP}{\mathcal{P}}
\newcommand{\calQ}{\mathcal{Q}}
\newcommand{\calR}{\mathcal{R}}
\newcommand{\calS}{\mathcal{S}}
\newcommand{\calT}{\mathcal{T}}
\newcommand{\calU}{\mathcal{U}}
\newcommand{\calV}{\mathcal{V}}
\newcommand{\calW}{\mathcal{W}}
\newcommand{\calX}{\mathcal{X}}
\newcommand{\calY}{\mathcal{Y}}
\newcommand{\calZ}{\mathcal{Z}}

\newcommand{\cc}{\mathfrak{c}}
\newcommand{\GG}{\mathfrak{G}}

%\DeclareMathOperator{\arg}{arg}
\DeclareMathOperator{\Aut}{Aut}
\DeclareMathOperator{\divv}{div}
\DeclareMathOperator{\Gal}{Gal}
\DeclareMathOperator{\GL}{GL}
\DeclareMathOperator{\HH}{H}
\DeclareMathOperator{\Hom}{Hom}
\DeclareMathOperator{\Inn}{Inn}
\DeclareMathOperator{\Ni}{Ni}
\DeclareMathOperator{\Spec}{Spec}
\DeclareMathOperator{\Stab}{Stab}

\newcommand{\Directsum}{\bigoplus}
\newcommand{\isom}{\simeq}
\newcommand{\isomto}{\stackrel{\sim}\to}
\newcommand{\surjects}{\twoheadrightarrow}

\newcommand{\arith}{\mathrm{arith}}
\newcommand{\et}{\mathrm{\acute{e}t}}
\newcommand{\geom}{\mathrm{geom}}

\newtheorem{theorem}{Theorem}[section]
\newtheorem{lemma}[theorem]{Lemma}
\newtheorem{corollary}[theorem]{Corollary}
\newtheorem{proposition}[theorem]{Proposition}

\theoremstyle{definition}
\newtheorem{definition}[theorem]{Definition}
\newtheorem{question}[theorem]{Question}
\newtheorem{conjecture}[theorem]{Conjecture}
\newtheorem{example}[theorem]{Example}
\newtheorem{examples}[theorem]{Examples}

\theoremstyle{remark}
\newtheorem{remark}[theorem]{Remark}
\newtheorem{remarks}[theorem]{Remarks}

\makeatletter
\g@addto@macro\bfseries{\boldmath} % This makes math in section titles bold.
\makeatother

\usepackage{microtype}  % This adjusts spacing between words so as to improve the probability of having line breaks in good places.

%%% \renewcommand{\baselinestretch}{3}   % for triple spacing

\usepackage[
%	draft,
%	colorlinks,
	pagebackref,
	pdfauthor={Xiaoyu Huang, Blake Jackson, Kyu-Hwan Lee, Bjorn Poonen, Rachel Pries, Shaowu Zhang}, 
	pdftitle={The Mathieu group M23 is a Galois group over Q},
]{hyperref}

\begin{document}

\title{The Mathieu group $M_{23}$ is a Galois group over $\Q$}
\subjclass[2020]{Primary 12F12; Secondary 11G32, 11R32, 12F10, 14H30, 14Q05}
\keywords{Inverse Galois problem, Mathieu group, sporadic group, Belyi map, rigidity method}

%SAVE
%12F12 Inverse Galois theory
%14H30 Coverings of curves, fundamental group
%14Q05 Computational aspects of algebraic curves

%11G32 Arithmetic aspects of dessins d’enfants
%11R32 Galois theory (global fields)
%12F10 Field extensions, Galois theory

\author[X. Huang]{Xiaoyu Huang}
\address{Department of Mathematics, Temple University, Philadelphia, PA 19122, USA} \email{xiaoyu.huang@temple.edu}
\urladdr{\url{https://sites.google.com/view/xiaoyuhuang/home}}

\author[B. Jackson]{Blake Jackson}
\address{Institute for Computer-Aided Reasoning in Mathematics, Carnegie Mellon University, Pittsburgh, PA 15213, USA}
\email{jackson@icarm.io}
\urladdr{\url{https://www.blakejacksonmath.com/}}

\author[K.-H. Lee]{Kyu-Hwan Lee}
\address{Department of Mathematics, University of Connecticut, Storrs, CT 06269, USA \hfill \break \indent Korea Institute for Advanced Study, Seoul 02455, Republic of Korea}
\email{khlee@math.uconn.edu}
\urladdr{\url{https://khlee-math.github.io/}}

\author[B. Poonen]{Bjorn Poonen}
\address{Department of Mathematics, Massachusetts Institute of Technology, Cambridge, MA 02139-4307, USA}
\email{poonen@math.mit.edu}
\urladdr{\url{http://math.mit.edu/~poonen/}}

\author[R. Pries]{Rachel Pries}
\address{Department of Mathematics, Colorado State University, Fort Collins, CO 80523, USA}
\email{pries@colostate.edu}
\urladdr{\url{http://math.colostate.edu/~pries/}}

\author[S. Zhang]{Shaowu Zhang}
\address{Department of Mathematics, California Institute of Technology, Pasadena, CA 91125, USA}
\email{szhang7@caltech.edu}
\urladdr{\url{https://www.shaowuzhang.com/}}

\thanks{X.H.\ was supported by an AMS--Simons Travel Grant. B.J.\ was supported by ICARM through NSF grant DMS-2425401. B.P.\ was supported in part by NSF grant DMS-2601551 and Simons Foundation grant \#402472.
R.P.\ was supported by NSF grant DMS-2200418.}

\date{August 8, 2026}

\begin{abstract}
Researchers studying the inverse Galois problem realized 25 of the 26 sporadic finite simple groups 
as Galois groups over $\Q$ during 1984--1989.
We complete this program by proving that the last remaining sporadic group, the Mathieu group $M_{23}$, 
occurs as a Galois group over $\Q$.
In fact, we produce an explicit degree $23$ polynomial with rational 
coefficients whose splitting field has Galois group $M_{23}$ over $\Q$. 
To accomplish this, we use a non-rigid triple of conjugacy classes of $M_{23}$ and compute Belyi maps to construct an explicit regular Galois extension of $\Q(t)$ with Galois group $M_{23}$.
Essential for our computation is the numerical Belyi map algorithm developed and implemented by Klug, Musty, Schiavone, Sijsling, and Voight, inspired by ideas of Hejhal and Stark.
\end{abstract}

\maketitle

%****************************************************************************
\section{Introduction}\label{S:introduction}

\subsection{The inverse Galois problem}
\label{S:inverse Galois problem}

The inverse Galois problem asks whether or not every finite group $G$ occurs as the Galois group of 
some finite Galois extension of $\Q$.
In particular, one can ask the question for finite simple groups, which fall into $18$ infinite families and $26$ ``sporadic'' groups.
During 1984--1989, $25$ of the $26$ sporadic groups were realized as Galois groups over $\Q$.\footnote{Strictly speaking, the baby monster $B$ was realized only later, since there was a calculation error in the original construction of \cite{Hunt1986}, according to \cite[p.~161]{Malle-Matzat2018}.}
But the last, the Mathieu group $M_{23}$, remained unrealized until now, despite many attempts; see Section~\ref{S:previous}.

\begin{theorem}
\label{T:over Q}
There exists a Galois extension of $\Q$ with Galois group $M_{23}$.
\end{theorem}

\begin{example}
\label{Ex:M23 over Q}
The splitting field of 
\begin{align*}
    f(x) \colonequals x^{23} &- 184 x^{21} - 1150 x^{20} + 26151 x^{19} + 18400 x^{18} - 1808490 x^{17} \\
      &+ 1545462 x^{16} + 67672923 x^{15} - 42732528 x^{14} - 1333395744 x^{13} \\
      &+ 290615166 x^{12} + 10550424369 x^{11} + 3700476348 x^{10} + 35123826654 x^9 \\
      &- 194398310718 x^8 - 1023887308293 x^7 + 3961650395556 x^6 \\
      &+ 1949980486716 x^5 - 28142323927002 x^4 + 53599151839311 x^3 \\
      &- 46185312415788 x^2 + 19169943578802 x - 3150159884154
\end{align*}
is a Galois extension of $\Q$ with Galois group $M_{23}$, unramified outside $\{2,3,23\}$. We certified these claims with Magma~\cite{Magma}.\footnote{Our code is available at \url{https://github.com/shaowuz/m23isgalois}.}
\end{example}

% (SAVE) Magma or PARI-GP notation: x^23 - 184*x^21 - 1150*x^20 + 26151*x^19 + 18400*x^18 - 1808490*x^17 + 1545462*x^16 + 67672923*x^15 - 42732528*x^14 - 1333395744*x^13 + 290615166*x^12 + 10550424369*x^11 + 3700476348*x^10 + 35123826654*x^9 - 194398310718*x^8 - 1023887308293*x^7 + 3961650395556*x^6 + 1949980486716*x^5 - 28142323927002*x^4 + 53599151839311*x^3 - 46185312415788*x^2 + 19169943578802*x - 3150159884154

% (SAVE) The discriminant of the degree $23$ extension is $2^{36} \cdot 3^{18} \cdot 23^{30}$.

\subsection{Regular Galois extensions}

A \defi{$G$-extension} of a field $k$ is a Galois field extension $L/k$, together with an isomorphism 
$\Gal(L/k) \isom G$.
A finite extension of $k(t)$ is \defi{regular} if it contains no nontrivial algebraic extension of $k$.
If $G \ne \{1\}$ and a regular $G$-extension of $\Q(t)$ exists,
then, by Hilbert's irreducibility theorem, 
specializing $t$ to different rational numbers yields infinitely many $G$-extensions of $\Q$.
Thus, to prove Theorem~\ref{T:over Q}, we prove the following stronger result. 

\begin{theorem}
\label{T:over Q(t)}
There exists a regular Galois extension $L/\Q(t)$ with Galois group $M_{23}$.
\end{theorem}

\subsection{Previous work}
\label{S:previous}

The rigidity method, originally developed in \cite{Shih1974,Fried1977,Belyi1979,Matzat1979,Matzat1984,Matzat1985a,Thompson1984},
was instrumental for constructing regular Galois extensions of $\Q(t)$.
Variants of this method, often involving significant new ideas such as the braid group method of \cite{Matzat1989,Matzat1991,Fried-Voelklein1991}, were used to construct regular $G$-extensions of $\Q(t)$ for all the sporadic groups $G$ except $M_{23}$ \cite{Thompson1984,Matzat1985a,Hoyden-Siedersleben1985, Hunt1986, Hoyden-Siedersleben-Matzat1986, Matzat-Zeh-Marschke1986, Matzat-Zeh-Marschke1987, Malle88,
Pahlings1988,Pahlings1989}; see \cite[II.9]{Malle-Matzat2018} for a detailed exposition.

Prior attempts to realize $M_{23}$ over $\Q$ resulted only in realizations over other fields:
\begin{itemize}
\item Hoyden-Siedersleben~\cite{Hoyden-Siedersleben1985} and H\"afner~\cite{Haefner1987} used regular $M_{24}$-extensions of $k(t)$ to construct regular $M_{23}$-extensions of $k(t)$ for $k=\Q(\sqrt{-23})$ and $k=\Q(\sqrt{-7})$, respectively.
\item Granboulan \cite{Granboulan1996} used another $M_{24}$-extension to construct a regular $M_{23}$-extension over $k(t)$ for any field $k$ over which a certain conic $C$ has a $k$-point.
Unfortunately, $C$ has no $\Q$-point.
\item Elkies \cite{Elkies2013} calculated the four degree~23 polynomials $P(x) \in \C[x]$ (up to affine equivalence) such that the Galois group of $P(x)-t$ over $\C(t)$ is $M_{23}$.
Each such $P(x)$ has coefficients in a degree~$4$ number field $F$, 
so one obtains a regular $M_{23}$-extension of $F(t)$.
\end{itemize}
For more information and explicit computation of Hurwitz spaces for $M_{23}$-covers, see \cite{Koenig2014}, \cite{Haefner-preprint}, and \cite{Seguin25}.

\subsection{Outline of our approach}

\begin{enumerate}[\upshape 1.]
\item The group $M_{23}$ has $17$ conjugacy classes, traditionally labeled by the order of a representative element followed by an upper case letter to distinguish classes of the same order.
The group $\GG_\Q \colonequals \Gal(\Qbar/\Q)$ acts on the set of conjugacy classes of $M_{23}$.
It fixes the class $2 \colonequals 2A$, and acts on $23A$ and $23B$ nontrivially through the quotient $\Gal(K/\Q)$, where $K \colonequals \Q(\sqrt{-23})$.
\item
Let $S$ be the finite \'etale $\Q$-scheme whose geometric points are $0,\sqrt{-23},-\sqrt{-23}$.
\item 
The Riemann existence theorem, together with a calculation in $M_{23}$, shows that there are exactly 
seven $M_{23}$-covers $Y_\C \to \PP^1_\C$ ramified only over $S$ with local monodromy $(2,23A,23B)$.
The quotient of each $Y_\C$ by the one-point stabilizer is a genus~$4$ curve $X_\C$.
\item
We compute these seven curves $X_\C$ numerically using the BelyiDB package of \cite{Musty-Schiavone-Sijsling-Voight2019},
and recognize coefficients as algebraic numbers to verify the output.
\item 
One might expect $\GG_\Q$ to act transitively on these seven.
Miraculously, instead $\GG_\Q$ has a fixed point! 
We do not have a conceptual explanation for this.
\item
The fixed point represents an $M_{23}$-cover $Y_\C \to \PP^1_\C$ with field of moduli $\Q$.
By \cite[Proposition 2.8(c)]{CoombesHarbater} %$G$-Galois covers
or
\cite[Proposition~3.1(b) and Corollary~3.2]{DebesDouai}, %mere covers
it descends to an $M_{23}$-cover $Y \to \PP^1_\Q$. 
\end{enumerate}

Details will be given in the subsequent sections.

\subsection{Conclusions}

Let $k$ be a number field.
Call a collection of $G$-extensions of $k$ \defi{mutually independent} if, for any $r \ge 1$, any $r$ distinct extensions from the collection generate a $G^r$-extension of $k$.

\begin{corollary}
Let $G$ be a sporadic finite simple group.
Let $k$ be any number field.
\begin{enumerate}[\upshape (a)]
\item \label{I:sporadic over Q(t)}
    There exists a regular $G$-extension of $k(t)$.
\item \label{I:sporadic over Q}
    There exists a $G$-extension of $k$.
\item \label{I:linearly disjoint}
    There exists a $G$-extension of $\Q$ linearly disjoint from $k$.
\item \label{I:mutually independent}
    There exists an infinite collection of mutually independent $G$-extensions of $k$.
\end{enumerate}
\end{corollary}

\begin{proof}
Part~\eqref{I:sporadic over Q(t)} over $\Q$ is the combination of Theorem~\ref{T:over Q(t)} and the work cited in Section~\ref{S:previous}.
Take a compositum to get \eqref{I:sporadic over Q(t)} over $k$.
For (\ref{I:sporadic over Q},\ref{I:linearly disjoint},\ref{I:mutually independent}) and other consequences of \eqref{I:sporadic over Q(t)}, see \cite[\S3.3--3.4]{Serre1992} and \cite[\S10.1]{Serre1997}.
\end{proof}

%****************************************************************************
\section{The rigidity method}

\subsection{\texorpdfstring{$G$}{G}-covers of the projective line}

Let $G$ be a nonabelian finite simple group, so $G$ has trivial center.
Let $\Inn(G)$ be the group of inner automorphisms of $G$, so $\Inn(G) \isom G$, and $\Inn(G)$ acts on $G$.
The set $\calC \colonequals \Inn(G) \backslash G$ is the set of conjugacy classes of $G$.
For $g \in G$, let $[g] \in \calC$ be its class.
Let $n \ge 1$ be such that $g^n=1$ for all $g \in G$.

Let $k$ be a field of characteristic~$0$.
For each $m \ge 1$, let $\mu_m = \{x \in \kbar^\times : x^m=1\}$.
Let $\GG_k = \Gal(\kbar/k)$.
Let $\epsilon \colon \GG_k \to \Aut \mu_n \isom (\Z/n\Z)^\times$ be the cyclotomic character describing the $\GG_k$-action on $n$th roots of unity.
A choice of generator $\zeta \in \mu_n$ determines a bijection $\Hom(\mu_n,G) \isom \Hom(\Z/n\Z,G) = G$.
To make this bijection \emph{$\GG_k$-equivariant}, we let $\sigma \in \GG_k$ act trivially on the $G$ in $\Hom(\mu_n,G)$ and as $g \mapsto g^{\epsilon(\sigma)^{-1}}$ on the $G$ on the right.
Taking $\Inn(G)$-orbits, the latter action induces a $\GG_k$-action on $\calC$.

A \defi{curve} over $k$ is a $1$-dimensional smooth projective geometrically connected scheme of finite type over $k$.
A \defi{$G$-cover of $\PP^1_k$} is a (ramified) Galois morphism of curves $Y \to \PP^1_k$ together with an isomorphism $\Gal(k(Y)/k(\PP^1)) \isom G$.
An \defi{isomorphism of $G$-covers} $Y \to \PP^1_k$ and 
$Y' \to \PP^1_k$ is a $G$-equivariant isomorphism $Y \to Y'$ over $\PP^1_k$.

\subsection{Local monodromy}

Let $D \subset \C$ be the open unit disk.
Let $D^\times = D -\{0\}$.
We have $\Z \isomto \pi_1(D^\times)$ sending $1$ to the standard counterclockwise generator.

Suppose that $Y_\C \to \PP^1_\C$ is a $G$-cover with branch locus $S$.
For each $s \in S$, we have an analytic open embedding $D \to \PP^1_\C$ sending $0$ to $s$, 
such that $s$ is the only point of $S$ in the image;
ignoring basepoints, this induces 
\[
\Z \isom \pi_1(D^\times) \to \pi_1(\PP^1(\C)-S) \surjects G
\]
sending $1$ to a generator of the local monodromy group (inertia group) at $s$.
Changing basepoints conjugates the generator by an element of $G$,
so we obtain a well-defined conjugacy class for each $s$, hence a map $\cc \colon S \to \calC$, 
which we call the \defi{local monodromy}.

Algebraically, if $u$ is an analytic uniformizer at $s \in \PP^1(\C)$, 
the completion of $\C(\PP^1)$ at $s$ is $\C((u))$, and we have analogous homomorphisms involving \'etale fundamental groups
\[
\Zhat(1) \isom \pi_1^{\et}(\Spec \C((u))) \to \pi_1^{\et}(\PP^1_\C-S) \surjects G,
\]
factoring through $\mu_n$, defined up to conjugation by elements of $G$.
Thus we get a canonical map
$\cc \colon S \to \Inn(G) \backslash \Hom(\mu_n,G) \isom \calC$, where the last isomorphism is induced by the choice $\zeta \colonequals e^{2\pi i/n}$.

If $Y_\C \to \PP^1_\C$ is the base change of a $G$-cover $Y \to \PP^1_\Q$,
then $S$ may be viewed as a finite \'etale subscheme of $\PP^1_\Q$,
and $\cc \colon S(\Qbar) \to \calC$ is $\GG_\Q$-equivariant.
This is the branch cycle argument of Fried.

\subsection{Rigidity}
\label{S:rigidity}

See \cite{Fried-Voelklein1991}, especially Corollary~1, for details on the theory in this section.
Now fix a finite \'etale subscheme $S \subset \PP^1_\Q$ and a $\GG_\Q$-equivariant map $\cc \colon S(\Qbar) \to \calC$. 
Write $S(\C) = \{s_1,\ldots,s_r\}$, and $C_i \colonequals \cc(s_i)$ for $i=1,\ldots,r$.


Let $\calH$ be the $0$-dimensional Hurwitz scheme over $\Q$ parameterizing $G$-covers of $\PP^1$ unramified outside $S$ with local monodromy $\cc$.
Because $G$ has trivial center, $\calH$ is a fine moduli space.
Define 
\begin{align*}
   \bar{\Sigma}_\cc &\colonequals \{(g_1,\ldots,g_r) \in C_1 \times \cdots \times C_r : g_1 \cdots g_r = 1 \} \\
   \Sigma_\cc &\colonequals \{(g_1,\ldots,g_r) \in \bar{\Sigma}_\cc : \textup{$g_1,\ldots,g_r$ generate $G$}\} \\
   \Ni_\cc &\colonequals \Inn(G) \backslash \Sigma_\cc;
\end{align*}
here $g \in G \isom \Inn(G)$ acts on $r$-tuples by simultaneous conjugation.
The set $\Ni_\cc$ is called the \defi{Nielsen class}.
By the Riemann existence theorem, $\calH(\C)$ is in bijection with $\Ni_\cc$. 

The size $\lvert \bar{\Sigma}_\cc \rvert$ can be computed from the character table of $G$; an inclusion--exclusion gives $\lvert \Sigma_\cc \rvert$; and dividing by $|G|$ gives $\lvert \Ni_\cc \rvert$, which depends only on the multiset $\{C_1,\ldots,C_r\}$; see \cite[\S7.2 and \S7.3]{Serre1992} for an exposition.

We restrict to the case $r=3$.
If $\lvert \Ni_\cc \rvert = 1$, then the triple $(C_1,C_2,C_3)$ is called \defi{rigid};
in that case, $\calH \isom \Spec \Q$ and the unique point corresponds to a $G$-cover $Y \to \PP^1_\Q$ with local monodromy $\cc$.

%****************************************************************************
\section{The Mathieu group \texorpdfstring{$M_{23}$}{M23}}

The Mathieu group $M_{23}$ is a finite simple group of order $10{,}200{,}960 = 2^7 \cdot 3^2 \cdot 5 \cdot 7 \cdot 11 \cdot 23$.
It is a $4$-transitive subgroup of $S_{23}$, generated by the permutations $g_1,g_2$ below.
It has 17 conjugacy classes:
\begin{center}
  1, \; 2, \; 3, \; 4, \; 5, \; 6, \; 7A, 7B, \; 8, \; 11A, 11B, \; 14A, 14B, \; 15A, 15B, \; 23A, 23B
\end{center}
(we drop the A in the label when there is a unique conjugacy class of elements of a particular order).
We calculate $\lvert \Ni_\cc \rvert$ for every $\GG_\Q$-stable $3$-element multiset in $\calC$,
using character tables as in Section~\ref{S:rigidity}.
The value of $\lvert \Ni_\cc \rvert$ is never $1$, so the rigidity method fails.

We next investigate the multiset for which $\lvert \Ni_\cc \rvert$ has smallest positive size.
This is $\{2,23A,23B\}$, with $\lvert \Ni_\cc \rvert = 7$.
In fact, $M_{23}$ is small enough that a computer can quickly list triples representing the elements of $\Ni_\cc$.
We fix the following such triple of permutations in $S_{23}$.\footnote{Our convention is that $S_{23}$ acts on the left on $\{1,\ldots,23\}$, but Magma has $S_{23}$ acting on the right.}
\begin{align*}
   g_1 &\colonequals (1,11)(2,23)(3,8)(4,16)(5,21)(7,20)(15,19)(18,22)  \\
   g_2 &\colonequals (1,2,11,10,16,9,6,3,23,19,20,14,21,17,4,8,22,5,18,15,13,7,12) \\
   g_3 &\colonequals (1,2,3,4,10,11,12,7,19,18,8,6,9,16,17,21,22,5,14,20,13,15,23). %(SAVE) this swapped g2 and g3 from Coco's code because we want a left action instead of a right action, which is what Magma uses.
\end{align*}

Let $S$ be the finite \'etale $\Q$-scheme whose geometric points are $0,\sqrt{-23},-\sqrt{-23}$.
Let $\cc \colon S(\Qbar) \to \calC$ send these points to classes 2, 23A, 23B, respectively.
This $\cc$ is $\GG_\Q$-equivariant.
We have $(g_1,g_2,g_3) \in \Sigma_\cc$, representing a class in $\Ni_\cc$.

Let $Y_\C \stackrel{\phi}\to \PP^1_\C$ be the corresponding $M_{23}$-cover.
Let $X_\C$ be the (variety) quotient of $Y_\C$ by $\Stab_{M_{23}}(1) \le M_{23}$,
so $\phi$ factors as $Y_\C \to X_\C \stackrel{t}\to \PP^1_\C$.
Then $X_\C \stackrel{t}\to \PP^1_\C$ is a \emph{non-Galois} degree~$23$ cover branched over $S(\C)$,
and its Galois closure is $Y_\C \to \PP^1_\C$.

\begin{lemma}
    The genus of $X_\C$ is $4$.
\end{lemma}

\begin{proof}
There are three ramified fibers of $t$, one with eight points of ramification index $2$ and seven unramified points, and two that are totally ramified. 
The degree of the ramification divisor of $t$ is $8(2-1) + 2(23-1)$, so the genus of $X_\C$ is $4$ by the Riemann--Hurwitz formula.
\end{proof}

Model the hyperbolic plane by the open unit disk $D$.  
Let $q$ be the standard coordinate on $D$.  
Let $d(z,w)$ be the metric on $D$.
Given $\theta \in (0,\pi/4)$, define
\[
    a \colonequals e^{-i \theta} \sqrt{\tan\left(\frac{\pi}{4}-\theta\right)}, \qquad b \colonequals 0, \qquad c \colonequals \sqrt{\cos 2\theta}.
\]

\begin{lemma}
  The points $a,b,c$ are the vertices of a hyperbolic triangle in $D$ with corresponding interior angles $\pi/2$, $\theta$, $\theta$. 
\end{lemma}

\begin{proof}
Note that $|a|^2 < 1$ and $|c|^2 < 1$, so $a,b,c \in D$.
Since $b=0$, we have $\angle b = \arg(c) - \arg(a) = \theta$.
The hyperbolic law of cosines gives $\cosh d(a,c) = \cot \theta$.
A formula for the cosines of the angles in terms of the side lengths of a hyperbolic triangle gives $\cos(\angle a) = 0$ and $\cos(\angle c) = \cos \theta$, so $\angle a = \pi/2$ and $\angle c = \theta$.
\end{proof}

Set $\theta \colonequals \pi/23$.
Let $\Delta$ be the triangle group generated by counterclockwise hyperbolic rotations by angles $2\pi/2$, $2\pi/23$, $2\pi/23$ around $a,b,c$: 
\[
\Delta \colonequals \langle \, \delta_a,\delta_b,\delta_c \mid \delta_a^2=\delta_b^{23}=\delta_c^{23}=\delta_a\delta_b\delta_c=1 \, \rangle.   
\]
Let $\psi \colon \Delta \to M_{23}$ be the surjection sending $\delta_a,\delta_b,\delta_c$ to $g_1,g_2,g_3$, respectively.
Let $\Gamma \colonequals \psi^{-1}(\Stab_{M_{23}}(1)) \le \Delta$.
Then $\Gamma \backslash D \to \Delta \backslash D$ is the desired cover 
$X_\C \stackrel{t}\to \PP^1_\C$ with monodromy given by $(g_1,g_2,g_3)$.
Let $a',b',c'$ be the images of $a,b,c$ under $D \surjects X_\C$.
The map $D \surjects X_\C$ is unramified at $b$ since the ramification indices of $b$ and $b'$ over $\PP^1$ are both $23$.
Thus $q$ corresponds to an analytic uniformizer at $b' \in X_\C$.

A basis of $\HH^0(X_\C,\Omega^1)$ pulls back to a $4$-tuple $(f_0 \, dq,\ldots,f_3 \, dq)$ for some $f_0,\ldots,f_3 \in \C[[q]]$.
Given $(g_1,g_2,g_3)$ and $n \ge 1$, the algorithm in \cite{Klug-Musty-Schiavone-Voight2014}, inspired by work of Hejhal~\cite{Hejhal1999} and Stark~\cite{Stark1984}, 
computes the first $n$ coefficients of each $f_i$ numerically.
The BelyiDB Magma package \cite{Musty-Schiavone-Sijsling-Voight2019} implements this algorithm.
We run it!
Below, we work numerically without controlling errors, so the intermediate steps of the calculation are not rigorous, but we return to rigor at the end by verifying algebraically that the claimed polynomials define $M_{23}$-extensions of $\Q$ and a regular $M_{23}$-extension of $\Q(t)$. 

We echelonize the basis.
This calculation shows that no nonzero $1$-form vanishes at $b'$ to order greater than $3$
(that is, $b'$ is not a Weierstrass point).
In other words, we obtain $f_i = q^i + O(q^4)$. 

\begin{lemma} 
The genus~$4$ curve $X_\C$ is non-hyperelliptic.
\end{lemma}

\begin{proof}
The dimension of the space of quadratic forms $P(x_0,\ldots,x_3)$ vanishing at $(f_0:\cdots:f_3)$ 
is $\binom{4-1}{2} = 3$ if $X_\C$ is hyperelliptic, and $\binom{4-2}{2} = 1$ otherwise \cite[Lemma~4.3]{Bommel-Costa-Poonen-Srinivasan2026}.
We compute that the dimension is $1$.
\end{proof}

The canonical model of a non-hyperelliptic curve of genus~$4$ is a complete intersection $P=Q=0$ in $\PP^3$, where $P$ and $Q$ are homogeneous of degrees $2$ and $3$, respectively.
To make $P$ unique, not just up to a scalar, we require that its leading coefficient be $1$ (with respect to some monomial order).
To make $Q$ unique, we require that $Q$ be orthogonal to $x_0 P,\ldots,x_3 P$ with respect to the Hermitian  pairing associated to the monomial basis and require $Q$ to have leading coefficient $1$.

Since $\lvert \Ni_\cc \rvert = 7$, the cover $t \colon X_\C \to \PP^1_\C$ associated to $(g_1,g_2,g_3)$ is defined over a number field $L$ of degree $\le 7$ over $\Q$.
We expanded the $1$-forms around the unique ramification point of $t \colon X_\C \to \PP^1_\C$ above $\sqrt{-23} \in \PP^1(L(\sqrt{-23}))$,
so the theory says that $P$ and $Q$ have coefficients in $L(\sqrt{-23})$.

Using the PSLQ algorithm \cite{PSLQ}, we recognize the coefficients of $P$ and $Q$ as algebraic numbers.
Miraculously, they lie in $K \colonequals \Q(\sqrt{-23})$! 
The exact $P$ and $Q$ define a canonically embedded genus~$4$ curve $X_K \subset \PP^3_K$.
Evaluating $(f_0:\cdots:f_3)$ at $a,b,c$ gives $a',b',c' \in X_K(\C) \subset \PP^3(\C)$
with $b' = (1:0:0:0)$, and PSLQ recognizes $c' \in X_K(K) \subset \PP^3(K)$. 

From now on, our calculations are algebraic.
We seek a morphism $t \colon X \to \PP^1_\Q$ which is ramified over $0$, $\sqrt{-23}$, $-\sqrt{-23}$, with ramification index $23$ above $\pm \sqrt{-23}$.
First we will construct a variant morphism $u$ ramified over $1,0,\infty$ instead.
A Riemann--Roch space computation produces $u \in K(X_K)$ such that $\divv(u) = 23b' - 23c'$.
We compute the zeros of $du$ other than $b'$ (the $8$ points with ramification index $2$)
and impose the condition that $u$ maps one of them to $1$, to obtain the correct scaling of $u$.
Let $\alpha \in \Aut(\PP^1_K)$ send $1,0,\infty$ to $0, \sqrt{-23}, -\sqrt{-23}$, respectively.
We compute $t = \alpha \circ u$.

We seek a second rational function $v \in K(X_K)$ such that $\Q(t,v)$ is the function field of a curve $X$ over $\Q$ whose base change is $X_K$.
Then the minimal polynomial of $v$ over $\Q(t)$ is the degree~$23$ polynomial whose splitting field is a regular $M_{23}$-extension of $\Q(t)$.

Let $\sigma$ be the generator of $\Gal(K/\Q)$, which acts also on $K(t)$, fixing $t$.
The $K$-vector space $L(5 b'-c')$ turns out to be $1$-dimensional, and a Riemann--Roch space computation produces a nonzero function $y$ in it.
We have $\deg y = 5$.
By linear algebra, we find $H(T,Y) \in K[T,Y]$ with $\deg_T H = 5$ and $\deg_Y H = 23$
such that $H(t,y)=0$.
We factor ${}^\sigma\! H(t,Y)$ over $K(X_K) = K(t,y)$ and find $y' \in K(X_K)$ such that $({}^\sigma\! H)(t, y')=0$.
This $y'$ is consistent with equaling ${}^\sigma \! y$ with respect to the intended $\Q$-structure on $X_K$.
Let $v=yy'$, which has polar divisor $4b'+4c'$ and is hence of degree $8$.
We next compute a polynomial $F(T,V) \in \Z[T,V] \subset \Q[T,V]$
such that $F(t,v)=0$, with $\deg_T F = 8$ and $\deg_V F = 23$;
to carry this out, we do linear algebra modulo small primes, and reconstruct coefficients using the Chinese remainder theorem.
Then $F(t,V) \in \Q(t)[V]$ is the desired degree~$23$ polynomial.
This polynomial is available on the GitHub site mentioned in the footnote at the end of Section~\ref{S:inverse Galois problem}.

\begin{remark}
In principle, we could also find a function $v$ of degree~$4$ instead of $8$,
which would reduce $\deg_T F$ to $4$.
To see this, let $K_X$ be a canonical divisor on the $\Q$-model $X$.
Then $K_X-(b'+c')$ is a divisor of degree $4$ defined over $\Q$.
Since $X$ is not hyperelliptic, $\dim L(b'+c')=1$,
so the Riemann--Roch theorem gives $\dim L(K_X-(b'+c'))=2$.
If $j_1,j_2$ is a $\Q$-basis, we can choose $v=j_1/j_2$.

In fact, degree~$4$ is best possible.
The only nonconstant rational maps $X_\C \to \PP^1_\C$ of degree $\le 3$ come from the rulings on the quadric $P=0$.
But the discriminant of $P$ is not a square in $K$,
so the rulings are not defined over $K$, let alone over $\Q$.
\end{remark}

Let $G_\geom$ and $G_\arith$ be the geometric and arithmetic monodromy groups of $F \in \Q(t)[V]$.
Thus $G_\geom \lhd G_\arith \le S_{23}$.

\begin{proposition} \label{PgeoM}
We have $G_\geom = M_{23}$. 
\end{proposition}

\begin{proof} 
The local monodromy above each of $\pm \sqrt{-23}$ is a $23$-cycle,
and above $0$ it has cycle type $1^7 2^8$.
The classification of transitive groups of degree $23$ then implies that $M_{23} \subseteq G_{\geom}$.

We calculate that the discriminant of $F \in \Q(t)[V]$ factors as $c t^8 (t^2+23)^{88} h_{84}^2$ for some nonzero $c \in \Z$ and degree~$84$ irreducible polynomial $h_{84} \in \Z[t]$.
(The presence of $h_{84}(t)$ reflects singularities of the plane model $F(T,V)=0$, not actual ramification of $X \to \PP^1_\Q$; we omit the proof since our argument does not rely on this.)
Moreover, the extension defined by $F$ is unramified above $\infty$.
Let $S'$ be the subscheme of $\PP^1_\Z$ defined by $t(t^2+23)h_{84}=0$.
We check that $31 \nmid c$, and $S'$ remains \'etale modulo $31$.
Also, $F \bmod 31$ defines a tame extension of $\F_{31}(t)$ since $31>23$.
Let $G_{\geom,31}$ and $G_{\arith,31}$ be the geometric and arithmetic monodromy groups of $F \bmod 31$.
The specialization map on tame fundamental groups of $\PP^1 - S'$ is an isomorphism (\cite[XIII.2]{SGA1}; see also \cite[Th\'eor\`eme~4.4]{Orgogozo-Vidal2000}), so $G_{\geom,31}=G_{\geom}$.
Using Magma, we certify that $G_{\arith,31} = M_{23}$.
Thus $M_{23} \subseteq G_{\geom} = G_{\geom,31} \subseteq G_{\arith,31} = M_{23}$,
so all are equal.
\end{proof}

\begin{corollary}
\label{C:G_arith}
We have $G_\arith=M_{23}$.
\end{corollary}

\begin{proof}
By Proposition~\ref{PgeoM}, $G_\geom = M_{23}$.
The group $G_\geom$ is normal in $G_\arith$,
and $M_{23}$ is not normal in any larger subgroup of $S_{23}$.
\end{proof}

By Proposition~\ref{PgeoM} and Corollary~\ref{C:G_arith},
since $G_\geom = G_\arith$, the splitting field of $F$ over $\Q(t)$ is a regular extension
with Galois group $M_{23}$.  This proves Theorem~\ref{T:over Q(t)}.

Finally, we specialize $t$ to rational numbers to obtain degree~$23$ polynomials over $\Q$ with Galois group $M_{23}$.  Then we use the PARI/GP routines \texttt{polredabs} and \texttt{polredbest} \cite{Cohen1991, PARI} to find lower-height monic polynomials in $\Z[x]$ defining the same degree~$23$ number field.
This is how we produced the polynomial in Example~\ref{Ex:M23 over Q}.

\begin{example}
\label{Ex:M23 over Q again}
Here is another specialization, the one of lowest height we found so far:
\begin{align*} 
f(x) \colonequals x^{23} &+ 46 x^{21} - 598 x^{20} + 1679 x^{19} - 21620 x^{18} + 127420 x^{17}\\ 
&- 361974 x^{16} + 2223732 x^{15} - 9392096 x^{14} + 17344116 x^{13}\\ 
&-71999476 x^{12} + 320807726 x^{11} - 436105484 x^{10} + 83587888 x^9\\ 
&-2463757240 x^8 + 9451874955 x^7 - 5728074376 x^6 - 26037806834 x^5\\ 
&+63691532334 x^4 - 67357061907 x^3 + 38754121124 x^2\\ 
&-11217790920 x + 1243077066. 
\end{align*}
Its splitting field is an $M_{23}$-extension of $\Q$ unramified outside $\{2,7,23\}$.
We certified this with Magma as well.
\end{example}


%****************************************************************************
\section*{Acknowledgments} 

We are very grateful to Michele Tarquini and Zhiyu Zhang for discussions when we were learning the theory and exploring approaches at the beginning of our project.
We began this research at the May 27--30, 2026 workshop ``AI and number theory'' at the American Institute of Mathematics in collaboration with the Institute for Computer-Aided Reasoning in Mathematics.
No text in this article was written by AI.
We used Claude Fable 5, Claude Opus 4.8, and ChatGPT 5.6 Sol for searching the literature, code generation, testing hypotheses, ruling out other approaches, devising computational strategies, checking our results, and proofreading our manuscript.
We thank Anthropic for providing each of us with a 3-month Claude Max subscription.
Our final results were verified in Magma and PARI/GP without the use of AI.

\bibliographystyle{amsalpha}
\bibliography{M23}



\end{document}
